% 全期間累積の変更履歴（課題3 用）
% revision_history/revision_history.tex から \input される．
%
% ライフサイクル:
%   - 提出のタイミングで，sections/changelog_current.tex の履歴行を
%     このファイルの末尾（\endhead と \bottomrule の間）にコピーする
%   - このファイルは「リセットしない」．全期間の履歴がここに蓄積される
%
% 書式:
%   YYYY-MM-DD & 学籍番号 : 氏名 & 対象（章・節） & 変更内容 \\
%
% 前提パッケージ（呼び出し側でロード済みであること）:
%   \usepackage{longtable}
%   \usepackage{booktabs}
%   \usepackage{array}

\begin{longtable}{p{2.2cm}p{3.5cm}p{4cm}p{5.5cm}}
\toprule
日付 & 担当 & 対象（章・節） & 変更内容 \\
\midrule
\endfirsthead

\toprule
日付 & 担当 & 対象（章・節） & 変更内容 \\
\midrule
\endhead

% --- 全期間の履歴行（提出時に changelog_current.tex から追記される） ---
2026-05-29 & 25G1099 : 長山 魁良 & 第3章 担当C詳細設計（オブジェクト設計・関数設計） & シリアル通信内容に音の長さを追加し，変数一覧・関数一覧を更新 \\
2026-05-29 & 24G1056 : 慶田 優 & 第1章 作業計画（WBS） & 担当AにProcessing GUI操作・シリアル通信設計を追加，担当Bに予備振り前のキャリブレーション状態・輝度エンコード制御の設計を追加 \\
2026-05-29 & 24G1056 : 慶田 優 & 第1章 作業計画（ガントチャート） & 機材購入遅延により，ハードウェア先行からソフトウェア設計先行（→ハードウェア→テスト）のフローへ全体を変更 \\

\bottomrule
\endfoot
\end{longtable}
